\section{Countermodels and non-substitution laws}
\label{sec:countermodels}

The following finite examples separate the conclusions proved above from
the estimates which remain open.

\subsection{Zero completion need not minimize variation}

Fix \(U,V\ge3\) and prescribe the primitive datum
\[
 k(u,v)=1
 \qquad
 ((u,v)=1,\ u\le U,\ v\le V).
\]
Its canonical zero completion is
\[
 K^0(u,v)=\one_{u\le U,v\le V}\one_{(u,v)=1}.
\]
The constant rectangle
\[
 F(u,v)=\one_{u\le U,v\le V}
\]
belongs to the same completion class.  By
\cref{thm:terminal-rectangle-expansion},
\[
 \cV_\square(F)=1.
\]
The zero completion is much rougher.  Put \(L=\min(U,V)\).  For every
\(2\le n\le L-1\),
\[
 K^0(n,n)=K^0(n+1,n+1)=0,
 \qquad
 K^0(n+1,n)=K^0(n,n+1)=1,
\]
and hence
\[
 \Delta_{12}K^0(n,n)=-2.
\]
Therefore
\begin{equation}
 \cV_\square(K^0)\ge2(L-2).
 \label{eq:primitive-mask-variation-lower}
\end{equation}
The completion \(F\) reduces the \(d=1\) variation, but
\(F^{[d]}\) is nonzero for many \(d\) and its diagonal tail need not
vanish.  This is the basic variation-tail trade-off.

\begin{remark}
The example does not show that the full cost \(\cC_D(F)\) is smaller than
\(\cC_\square(K^0)\) at every scale.  It shows only that zero completion is not
variation-minimal, which is the precise claim needed to prevent a false
optimization statement.
\end{remark}

\subsection{Completion cost can be inflated arbitrarily}

Take the zero primitive datum \(k=0\), a prime \(p\le\min(U,V)\), and
\[
 F_A=Ae_{p,p}.
\]
The point \((p,p)\) is nonprimitive, so \(F_A\in\Ext_{U,V}(0)\) and the
physical functional is zero.  For \(D=1\), the exact tail contains the
term \(d=p,x=y=1\), giving
\[
 \cT_1(F_A)\ge |A|.
\]
Thus
\[
 \cC_1(F_A)\ge|A|
\]
can be arbitrarily large while the physical datum and physical sum remain
zero.  A large cost for an arbitrary analytic completion is not a physical
obstruction.

\subsection{Primitive support and tail freedom do not cancel}

Let \((u_0,v_0)=1\) and
\[
 K^0=ce_{u_0,v_0}.
\]
It is primitive-supported and has zero diagonal tail, but
\[
 \left|
 \sum_{u,v}\mu(u)\mu(v)K^0(u,v)
 \right|
 =
 |c|\,|\mu(u_0)\mu(v_0)|.
\]
When \(u_0,v_0\) are squarefree this equals the full \(\ell^1\)-mass
\(|c|\).  Tail freedom is therefore an exact bookkeeping improvement, not
an arithmetic cancellation theorem.

\subsection{Slice and anchor cancellation are different}

In the ambient finite-kernel class, two formal \(q\)-slices \(F\) and
\(-F\) have zero post-\(q\)-aggregation kernel although the sum of the two
separate costs is positive.  Likewise, an abstract anchor-compatible model
may have
\[
 \varepsilon_{\mathfrak a_1}K_{\mathfrak a_1}^{0}=F,
 \qquad
 \varepsilon_{\mathfrak a_2}K_{\mathfrak a_2}^{0}=-F.
\]
Then \(K_{\rm tr}^{0}=0\) and the scalar trace certificate vanishes, while
the anchor-labelled \(\ell^1\) cost is twice the cost of \(F\).
Consequently:
\begin{align*}
 \text{separate \(q\)-cost}
 &\not\equiv
 \text{post-\(q\)-aggregation cost},\\
 \text{scalar trace cancellation}
 &\not\equiv
 \text{anchorwise or row-local control}.
\end{align*}

\subsection{A vanishing mixed carrier does not control the erasure operator}

The abstract example
\[
 R=0,
 \qquad
 E=
 \begin{pmatrix}
 0&1\\
 1&0
 \end{pmatrix}
\]
has
\[
 \tr(ER)=0,
 \qquad
 \|E\|=1.
\]
In this abstract model the mixed carrier \(ER\) vanishes identically,
while the erasure operator is nonzero.  No scalar completion theorem
can remove the need for a genuine two-corner, row-local, or complete-moment
operator carrier.

\subsection{Quantifier-adaptive completion is invalid}

For each fixed primitive datum, finite-dimensional optimization can choose
a completion with small cost.  A coefficient-uniform operator theorem
requires a single declared completion rule valid for the entire physical
coefficient class.  Choosing a different completion after the maximizing
coefficient vector or row is known reverses the max--min order.  The
pointwise infimum in \eqref{eq:completion-gauge} is therefore an analytic
benchmark; it is not automatically a uniform L2 theorem.
